\section{问题重述}
\subsection{背景与任务}
通用神经网络处理器以多个 AI 核执行计算图。每核具有矩阵计算、向量计算与数据搬运等独立执行通道，但核内存储有限，多个核心共享核外 DDR 带宽。原始计算图为带张量依赖的有向无环图；切分粒度决定可同时运行的工作量，也改变输入重复读取、跨子图数据搬运和缓存溢出。题目要求输出每个非 \texttt{COPY\_IN}/\texttt{COPY\_OUT} 算子的子图编号，以及各核心上的子图顺序，再由给定评估程序确定总执行周期和额外搬运量\cite{problem}。

三个问题的决策形式相近，运行语义逐步改变。问题一的场景 A 规定一个子图就是一个 Task，子图边界上的数据总经 DDR 中转；即使两个 Task 在同一核上连续执行，也不保留私有缓存。问题二的场景 B 把同核各子图纳入一个 Task，允许核内驻留，但跨核数据需等待同步释放。问题三在场景 B 上增加全核共享的只读 FIFO L2，读请求可由 DDR 或 L2 服务。因而，问题二和三不能直接沿用场景 A 的通信成本公式，也不能把问题三的缓存命中率当作独立可选变量。

\subsection{所需回答与指标}
对每张图和每个核数 $k=1,\ldots,5$，本文构造合法的切图、核心分配与核上顺序。首要评价量为官方仿真得到的完成时间 $T_{i,k}$（cycles），同时报告逻辑额外数据搬运量 $D_{i,k}$（bytes）。设 $T_{i,1}^{0}$ 是题目规定的整图单核核内启发式结果，则前两问的逐图加速比与平均加速比定义为
\begin{equation}\label{eq:speedup}
S_{i,k}=\frac{T_{i,1}^{0}}{T_{i,k}},\qquad
\overline S_k=\frac1{100}\sum_{i=1}^{100}S_{i,k}.
\end{equation}
这是一百个比值的算术平均，不能换成总单核周期除以总多核周期。问题三另比较相同图、相同核数下的无 L2 与有 L2 周期，单核缓存收益由真实仿真决定，不人为设为 $1$。

\subsection{研究范围}
全部实验使用题目给定的 $100$ 张计算图和固定硬件参数。本文不改变图的算子、张量或官方核内启发式；搜索仅改变题目允许的划分、映射和子图顺序。求解器运行秒数与被评估方案的执行周期是不同量，前者反映求解效率，后者是模型的主要目标。各问的逐例结果列于附录，所有整体结论均以同一百图为统计总体。
